% Part: first-order-logic
% Chapter: syntax-and-semantics
% Section: subformulas

\documentclass[../../../include/open-logic-section]{subfiles}

\begin{document}

\olfileid{fol}{syn}{sbf}

\olsection{\printtoken{P}{subformula}}

\begin{explain}
Het is vaak nuttig om te spreken over de !!{formula}s waaruit een
gegeven !!{formula} is ``opgebouwd''. We noemen deze haar
\emph{!!{subformula}s}. Iedere !!{formula} geldt als !!a{subformula}
van zichzelf; een subformule van $!A$ die niet $!A$ zelf is, heet een
\emph{echte !!{subformula}}.
\end{explain}

\begin{defn}[Onmiddellijk !!^{subformula}]
Als $!A$ !!a{formula} is, worden de \emph{onmiddellijke
!!{subformula}s} van $!A$ als volgt inductief gedefinieerd:
\begin{enumerate}
\item Atomaire !!{formula}s hebben geen onmiddellijke !!{subformula}s.

\tagitem{prvNot}{\indcase{!A}{\lnot !B}{De enige onmiddellijke
    !!{subformula} van $\indfrm$ is~$!B$.}}{}

\item \indcase{!A}{(!B \ast !C)}{De onmiddellijke !!{subformula}s van
  $\indfrm$ zijn $!B$ en $!C$ ($\ast$ is een willekeurig tweeplaatsig
  connectief).}

\tagitem{prvAll}{\indcase{!A}{\lforall[x][!B]}{De enige onmiddellijke
    !!{subformula} van $\indfrm$ is~$!B$.}}{}

\tagitem{prvEx}{\indcase{!A}{\lexists[x][!B]}{De enige onmiddellijke
    !!{subformula} van $\indfrm$ is~$!B$.}}{}
\end{enumerate}
\end{defn}

\begin{defn}[Echt !!^{subformula}]
Als $!A$ !!a{formula} is, worden de \emph{echte !!{subformula}s} van
$!A$ als volgt recursief gedefinieerd:
\begin{enumerate}
\item Atomaire !!{formula}s hebben geen echte !!{subformula}s.

\tagitem{prvNot}{\indcase{!A}{\lnot !B}{De echte !!{subformula}s van
    $\indfrm$ zijn~$!B$ samen met alle echte !!{subformula}s
    van~$!B$.}}{}

\item \indcase{!A}{(!B \ast !C)}{De echte !!{subformula}s van
  $\indfrm$ zijn $!B$ en $!C$, samen met alle echte !!{subformula}s
  van $!B$ en die van~$!C$.}

\tagitem{prvAll}{\indcase{!A}{\lforall[x][!B]}{De echte
    !!{subformula}s van $\indfrm$ zijn~$!B$ samen met alle echte
    !!{subformula}s van~$!B$.}}{}

\tagitem{prvEx}{\indcase{!A}{\lexists[x][!B]}{De echte
    !!{subformula}s van $\indfrm$ zijn~$!B$ samen met alle echte
    !!{subformula}s van~$!B$.}}{}
\end{enumerate}
\end{defn}

\begin{defn}[!!^{subformula}]
De !!{subformula}s van $!A$ zijn $!A$ zelf samen met al haar echte
!!{subformula}s.
\end{defn}

\begin{explain}
Let op het subtiele verschil tussen de definities van onmiddellijke
!!{subformula}s en echte !!{subformula}s. In het eerste geval hebben
we voor iedere mogelijke vorm van een formule~$!A$ rechtstreeks de
onmiddellijke !!{subformula}s van~$!A$ gedefinieerd. Dit is een
expliciete definitie door gevallen; de gevallen weerspiegelen de
inductieve definitie van de verzameling !!{formula}s. In het tweede
geval weerspiegelen we eveneens de definitie van de verzameling van
alle !!{formula}s, maar nemen we naast de kleinere !!{formula}s $!B$,
$!C$ zelf ook telkens hun echte !!{subformula}s op. Daardoor is de
definitie \emph{recursief}. In het algemeen is een definitie van een
functie op een inductief gedefinieerde verzameling (hier de
!!{formula}s) recursief als de gevallen in de functiedefinitie de
functie zelf gebruiken. Voor een welgedefinieerde functie mogen we
daarbij alleen functiewaarden gebruiken voor argumenten die ``vóór''
het argument komen dat we definiëren. Bij de definitie van ``echt
!!{subformula}'' voor $(!B \ast !C)$ gebruiken we dus alleen de echte
!!{subformula}s van de ``eerdere'' !!{formula}s $!B$ en $!C$.
\end{explain}

\begin{prop}
\ollabel{prop:subfrm-trans}
Stel dat $!B$ een subformule van $!A$ is en $!C$ een subformule van
$!B$. Dan is $!C$ een subformule van $!A$. Met andere woorden: de
subformulerelatie is transitief.
\end{prop}

\begin{prob}
Bewijs \olref[fol][syn][sbf]{prop:subfrm-trans}.
\end{prob}

\begin{prop}
\ollabel{prop:count-subfrms}
Stel dat $!A$ een formule met $n$ connectieven en kwantoren is. Dan
heeft $!A$ ten hoogste $2n+1$ subformules.
\end{prop}

\begin{prob}
Bewijs \olref[fol][syn][sbf]{prop:count-subfrms}.
\end{prob}
\end{document}
